\subsection{Support reductions and the information they retain}

Several support constructions require only elementary algebra. Stating
their precise targets prevents them from being read as reconstructions
of the original arithmetic data.

\begin{proposition}[A multiplicative support quotient]
\label{prop:rees-support-quotient}
For a unital monoid $M$, the Rees quotient of $M\times\B$ by its absorbing
ideal $M\times\{0\}$ is $M$ with one new multiplicative absorber.
\end{proposition}
\begin{proof}
Collapse all $(m,0)$ to a single point $*$ and retain the classes of
$(m,1)$ separately. The product of two retained classes is $(mn,1)$;
any product involving the collapsed ideal is in that ideal. These are
exactly the asserted operations, and $(1,1)$ is the unit. This is a
statement about multiplicative monoids. No unspecified additive
operation or semiring isomorphism follows from it.
\end{proof}

\begin{proposition}[Initial data do not reconstruct a valued element]
\label{prop:initial-data}
Let $K$ be a valued field, with associated graded ring
$\operatorname{gr}_vK$. On $K_0$, the new zero and the old zero remain
distinct. For an old nonzero $x$, its valuation and initial class
$\operatorname{in}_v(x)\in(\operatorname{gr}_vK)_{v(x)}$ specify its first
graded term, but need not determine $x$. If a Laurent-series coordinate
is fixed, the initial class is equivalently the leading coefficient
together with its exponent.
\end{proposition}
\begin{proof}
For $x\ne0$, the initial class is its image in
$\{y:v(y)\ge v(x)\}/\{y:v(y)>v(x)\}$. Thus two elements of the same
valuation have the same initial class whenever their difference has
larger valuation. In $k((t))$, the distinct elements $t$ and $t+t^2$
have valuation $1$ and initial class $t$. They also have the same
occupied support in $K_0$. The two zero states are separately specified
by construction. Consequently a datum consisting of support, valuation,
and leading coefficient is initial data, not a complete invariant of
a Laurent series or any larger valued-field realization containing
this example.
\end{proof}

For a finite witness set $W$, its cardinality can be encoded inside
$\mathbb Z_0$ by the sum $c(W)=\sum_{w\in W}1$: this is the new zero for
$W=\varnothing$ and the old positive integer $|W|$ otherwise. Hence its
Boolean support is exactly the truth value of $W\ne\varnothing$.
Disjoint unions and Cartesian products give addition and multiplication
of these counts, including empty sets. Applied to any explicitly
specified finite Diophantine witness set, this records solvability
exactly, but does not by itself prove that a witness exists. It uses the
natural copy of $\mathbb N$ in $\mathbb Z_0$ whose zero is the new zero;
it makes no assertion that the larger zero-adjoined set $\mathbb N_0$
is an inverse semiring.
